\scanpage{40}
One concludes from the first group
\[
\varphi_*\Lambda_Y\varphi^*
=\underbrace{\varphi_*\varphi^*\Lambda_X}_{\substack{L_X\Lambda_X\\\Vert\\\operatorname{id}_X-\sum_{0\leq i\leq2n}p_i^X}}
\left(\operatorname{id}_X-\sum_{n+1\leq j\leq2n}p_X^j\right)
=\underbrace{\operatorname{id}_X-\sum_{0\leq j\leq2n}p_X^j}_{q}.
\]
[More generally
\[
\varphi_*\Lambda_Y^m\varphi^*
=\deleted{\sum}\left(\operatorname{id}-\sum_{0\leq i\leq n}p_X^i\right)\Lambda_X^{m-1}
\left(\operatorname{id}-\sum_{n+1\leq j\leq2n}p_X^i\right)
=q\Lambda_X^{m-1}q
\]
where $q=\operatorname{id}-\displaystyle\sum_{0\leq i\leq2n}p_X^i$.
\figure{\begin{tikzpicture}[x=1cm,y=1cm]\draw (0,.3)..controls(1,.5)and(2,.5)..(2.2,.15)..controls(2.3,-.1)and(2.25,-.3)..(2.6,-.3)--(7,-.4);\end{tikzpicture}}

Also:
\[
\varphi^*\Lambda_X\varphi_*=\operatorname{id}_Y-\underbrace{p_E}_{\text{projection onto vanishing part }E^{n-1}}
\]
whence, multiplying by $\Lambda_Y$ on the left, since $\Lambda_Yp_E=0$
\[
\Lambda_Y=\underbrace{\Lambda_Y\varphi^*}\Lambda_X\varphi_*
=\underbrace{\varphi^*\Lambda_X(\operatorname{id}_X-\sum_{n+1\leq j\leq2n}\overbrace{p_X^j)\Lambda_X}^{\substack{0\\\Vert}}}\varphi_*
\]
thus
\[\boxed{\Lambda_Y=\varphi^*\Lambda_X^2\varphi_*}\]
Thus $C(X)$ \deleted{\unsure{is equivalent to}}\inserted{implies} $C(Y)$.

Let $p_X^{i,j}$ be the projections onto $L_X^jP^i(X)$ $(i+j\leq n,\ i,j\geq0)$.

Then
\[
\varphi_*p_Y^{i,j}\inserted{L}\varphi^*=p_X^{i,j+1}
\qquad\text{for }i+j\unsure{<|\leq}n-1,\quad i,j\geq0\quad\deleted{\unsure{algebraic|arbitrary}}
\]
\margin{$L?$}
thus the condition $D(Y)$ implies that the $p_X^{i,j}$ are algebraic for $\unsure{\forall}\ i+j<n$, $i\geq0$ and \uline{$j\geq1$}, i.e. the only ones \unsure{still pending} are the \uline{$p_X^i$ for $0\leq i\leq2n$}. But
\mspage{39}
